
% 3. Environnement unique et superposition directe
\begin{tikzpicture}
  % On dessine les deux éléments sur le MÊME repère absolu
  %\clip[](-3.6, 3.4) rectangle (2.8 , -4.2)  ;

  \node[transform shape, scale=1] at (0,0) {
    \begin{tikzpicture}


        \def\a{0.04} \def\b{0.08}
        \def\h{20} \def\vOx{10} \def\g{9.81}

        %\draw[smooth, line width=0.01ex, color=colorFour, domain=-0.050:2.050, samples=30] plot (\x, \h - {0.5*\g/\vOx/\vOx*\x*\x});
        \draw[smooth, line width=0.1cm, color=colorFour, domain=-0.050:20.050, samples=30] plot (\x, {- 0.5 * \g / (\vOx * \vOx) * \x * \x});
        \begin{scope}[shift={(-0.0,0.0)}, xscale=1,yscale=1]
            % Grille principale
            \draw[opacity=0.2,very thin,line width=0.05cm,color=colorThree](-0.5,0.5)grid[step=1.0](20.5,-20.5);
            % Sous-grille
            \draw[opacity=0.2,very thin,line width=0.02cm,color=colorThree](-0.45,0.45)grid[step=0.5](20.45,-20.45);

            % Graduations X
            \foreach \x in {-0.5,0,0.5,1,...,20, 20.5} {
                \draw[line width=0.1cm,color=colorOne](\x,0-0.2)--(\x,0+0.2);
            }
            \foreach \x in {1,2,3,...,20} {
                \draw[line width=0.1cm,color=colorOne](\x,0-0.5)--(\x,0+0.5)
                node[shift={(0,0.5)},font=\fontsize{0.5cm}{0.5cm}\bfseries\selectfont]{\x};
            }
            % Graduations Y
             \foreach \y in {-0.5,-1,...,-20.5} {
                \draw[line width=0.1cm,color=colorOne,](0-0.2,\y)--(0+0.2,\y);
            }
            \foreach \y in {-1,-2,-3,...,-20} {
                \draw[line width=0.1cm,color=colorOne,](0-0.5,\y)--(0+0.5,\y)
                node[pos = 0 , left , shift={(-0.15,0)},font=\fontsize{0.5cm}{0.5cm}\bfseries\selectfont]{\y};
            }

            % 2. Axe des abscisses (x)
            \draw[line width=0.1cm,color=colorOne,->,
            >={Computer Modern Rightarrow[length=0.4cm, width=0.6cm]}
            ] (-1.0,0) -- (21.0,0) 
                node[pos=0.95, xshift=-0.00cm,yshift=-1.0cm, color=colorOne, font=\fontsize{0.5cm}{0.5cm}\bfseries\selectfont] {$\bm{x~(m)}$};

            % 3. Axe des ordonnées (S(x))
            \draw[line width=0.1cm,color=colorOne,->,
            >={Computer Modern Rightarrow[length=0.4cm, width=0.6cm]}
            ] (0,-21.1) -- (0,1.0) 
                node[pos=0.99, xshift=-1.2cm,yshift=-0.0cm, color=colorOne, font=\fontsize{0.5cm}{0.5cm}\bfseries\selectfont] {$\bm{y~(m)}$}
                ;

            
            \draw[shift={(12.00,-2.00)},line width=0.1cm,color=colorFour,<->,
            >={Computer Modern Rightarrow[length=0.4cm, width=0.6cm]}
            ] (0,0) -- (5,0) 
                node[pos=0.5, xshift=-0.0cm,yshift=0.5cm, color=colorFour, font=\fontsize{0.5cm}{0.5cm}\bfseries\selectfont] {$\bm{5 {~} m}$}
                ;

            \draw[shift={(12,-4)},line width=0.1cm,color=colorSix,<->,
            >={Computer Modern Rightarrow[length=0.4cm, width=0.6cm]}
            ] (0,0) -- (5,0) 
                node[pos=0.5, xshift=-0.0cm,yshift=0.5cm, color=colorSix, font=\fontsize{0.5cm}{0.5cm}\bfseries\selectfont] {$\bm{10 {~} m\cdot s^{-1}}$}%{$\bm{5\times10^{-2}} \mkern-20mu \mathrm{m\cdot s^{-1}}$}
                ;

            \draw[shift={(12,-6)},line width=0.1cm,color=colorFive,<->,
            >={Computer Modern Rightarrow[length=0.4cm, width=0.6cm]}
            ] (0,0) -- (5,0) 
                node[pos=0.5, xshift=-0.0cm,yshift=0.5cm, color=colorFive, font=\fontsize{0.5cm}{0.5cm}\bfseries\selectfont] {$\bm{9.81 {~} m\cdot s^{-2}}$}
                ;

        \end{scope}

        

        \draw[shift={(0,0)},opacity = 0.9 , line width=0.2cm,color=colorSix,->,
            >={Computer Modern Rightarrow[length=0.4cm, width=0.6cm]}
            ] (0,0) -- (5,0) 
                node[pos=0.5, opacity = 0.9 , xshift=-0.0cm,yshift=0.5cm, color=colorSix, font=\fontsize{0.5cm}{0.5cm}\bfseries\selectfont] {$\bm{\vec{v}_0 = 10 {~} m\cdot s^{-1} \vec{u}_x$}}
                ;

        \draw[shift={(3,-4)},opacity = 0.9 , line width=0.2cm,color=colorFive,->,
            >={Computer Modern Rightarrow[length=0.4cm, width=0.6cm]}
            ] (0,0) -- (0,-5) 
                node[pos=0.5, above , opacity = 0.9 , xshift=-0.0cm,yshift=0.0cm, color=colorFive, font=\fontsize{0.5cm}{0.5cm}\bfseries\selectfont , rotate = 90] {$\bm{\vec{g} = -9.81 {~} m\cdot s^{-2} \vec{u}_y$}}
                ;

        \draw [decoration={snake, amplitude=0.4cm, segment length=1.6cm},]
            (-1,-20) edge[ decorate, color=colorFour , line width=0.1cm,] (21,-20)
        ;
        
        \def\dt{0.1}

        \def\listeLabels{
            0/////{\vec{a}_0},
            1////{}/,
            2////{\vec{v}_2}/{\vec{a}_2},
            3/////,
            4/////{\vec{a}_4},
            5////{\vec{v}_5}/,
            6/////{\vec{a}_6},
            7/////,
            8/////,
            9/////,
            10/////,
            11/////,
            12/////,
        }

        % \foreach \n/\labelBottom/\labelTop/\LabelvecPos/\LabelvecVit/\LabelvecAcc in \listeLabels {%

        %     \pgfmathsetmacro{\t}{\n*\dt}
        %     \pgfmathsetmacro{\x}{\a*\t}
        %     \pgfmathsetmacro{\y}{-\b*\t*\t}

        %     \pgfmathsetmacro{\vx}{\a}
        %     \pgfmathsetmacro{\vy}{-2*\b*\t}

        %     \pgfmathsetmacro{\ax}{0}
        %     \pgfmathsetmacro{\ay}{-2*\b}

    
        %     % Point
        %     \def\taillecrois{0.0015}
        %     \fill[opacity = 0.5, color=colorFour] (\x,\y) circle (\taillecrois);
            
        %     \draw[opacity = 1,line width=0.001cm,shift={(0,0)}]
        %         (\x-\taillecrois,\y-\taillecrois)edge[shift={(0,0)}, opacity = 1 ,  color=colorThree]node[pos=0.5,shift={(0.0cm,0.0cm)},color=colorOne,font=\fontsize{0.1cm}{0.04cm}\bfseries\selectfont, rotate = 0]{$ $}(\x+\taillecrois,\y+\taillecrois)
        %         (\x-\taillecrois,\y+\taillecrois)edge[shift={(0,0)}, opacity = 1 ,  color=colorThree]node[pos=0.5,shift={(0.0cm,0.0cm)},color=colorOne,font=\fontsize{0.1cm}{0.04cm}\bfseries\selectfont, rotate = 0]{$ $}(\x+\taillecrois,\y-\taillecrois)
        %     ;

        %     % Label inférieur
        %     \if\relax\detokenize{\labelBottom}\relax
        %     \else
        %         \node[color=colorThree,shift={(0.01,0.01)},font=\fontsize{0.005cm}{0.005cm}\bfseries\selectfont] at (\x,\y) {$\bm{\labelBottom}$};
        %     \fi

        %     % Label supérieur
        %     \if\relax\detokenize{\labelTop}\relax
        %     \else
        %         \node[color=colorThree,shift={(0.02,0.0)},font=\fontsize{0.005cm}{0.005cm}\bfseries\selectfont] at (\x,\y) {$\bm{\labelTop}$};
        %     \fi

        %     \if\relax\detokenize{\LabelvecPos}\relax
        %     \else
        %     \fi

        %     \if\relax\expandafter\detokenize\expandafter{\LabelvecVit}\relax
        %     \else
        %         \draw[line width=0.001cm]
        %             (\x,\y) edge[dash pattern=on 0.002cm off 0.001cm on 0.002cm off 0.001cm, opacity=0.5, color=colorSix!50!colorTwo, ->,
        %                 >={Computer Modern Rightarrow[length=0.004cm, width=0.006cm]}]
        %                 node[pos=0.5, opacity=0.5,shift={(0.0,0.005)}, color=colorSix!50!colorTwo,
        %                     font=\fontsize{0.005cm}{0.005cm}\bfseries\selectfont]{$\bm{\vec{v}_x}$}
        %                 (\x+\vx,\y)
        %             (\x+\vx,\y) edge[dash pattern=on 0.002cm off 0.001cm on 0.002cm off 0.001cm, opacity=0.5, color=colorSix!50!colorTwo, ->,
        %                 >={Computer Modern Rightarrow[length=0.004cm, width=0.006cm]}]
        %                 node[pos=0.6, opacity=0.5, shift={(0.005,0.0)},color=colorSix!50!colorTwo,
        %                     font=\fontsize{0.005cm}{0.005cm}\bfseries\selectfont,
        %                     rotate=-90
        %                     ]{$\bm{\vec{v}_y}$}
        %                 (\x+\vx,\y+\vy)
        %             (\x,\y) edge[opacity=1, color=colorSix, ->,
        %                 >={Computer Modern Rightarrow[length=0.004cm, width=0.006cm]}]
        %                 node[pos=1.0, color=colorSix, shift={(0.000,0.005)},
        %                     font=\fontsize{0.005cm}{0.005cm}\bfseries\selectfont]{$\bm{\LabelvecVit}$}
        %                 (\x+\vx,\y+\vy);
        %     \fi

        %     \if\relax\expandafter\detokenize\expandafter{\LabelvecAcc}\relax
        %     \else
        %         \draw[line width=0.001cm]
        %             (\x,\y) edge[opacity=1, color=colorFive, ->,
        %                 >={Computer Modern Rightarrow[length=0.004cm, width=0.006cm]}]
        %                 node[pos=1.0 , color=colorFive, shift={(0.01,-0.005)},
        %                     font=\fontsize{0.005cm}{0.005cm}\bfseries\selectfont]{$\bm{\LabelvecAcc}$}
        %                 (\x+\ax,\y+\ay/2);
        %     \fi

        %     % \ifstrempty{\LabelvecVit}{}{%
        %     %     \draw[line width=0.001cm]
        %     %         (\x,\y) edge[dashed, opacity=0.05, color=colorOne, ->,
        %     %             >={Computer Modern Rightarrow[length=0.004cm, width=0.006cm]}]
        %     %             node[pos=1.1, opacity=0.05, color=colorOne,
        %     %                 font=\fontsize{0.005cm}{0.005cm}\bfseries\selectfont]{$\bm{\vec{v}_x}$}
        %     %             (\x+\vx,\y)
        %     %         (\x+\vx,\y) edge[dashed, opacity=0.05, color=colorOne, ->,
        %     %             >={Computer Modern Rightarrow[length=0.004cm, width=0.006cm]}]
        %     %             node[pos=1.1, opacity=0.05, color=colorOne,
        %     %                 font=\fontsize{0.005cm}{0.005cm}\bfseries\selectfont]{$\bm{\vec{v}_y}$}
        %     %             (\x+\vx,\y+\vy)
        %     %         (\x,\y) edge[opacity=1, color=colorOne, ->,
        %     %             >={Computer Modern Rightarrow[length=0.004cm, width=0.006cm]}]
        %     %             node[pos=1.1, color=colorOne,
        %     %                 font=\fontsize{0.005cm}{0.005cm}\bfseries\selectfont]{$\bm{\LabelvecVit}$}
        %     %             (\x+\vx,\y+\vy);
        %     % }

            
        % }

        % 4. Courbe de la parabole y = x^2 (domaine restreint pour rester dans le cadre)
        %\draw[smooth, line width=1.0ex, color=colorFour, domain=-2.1:2.1, samples=100] plot (\x, {1-2*\x*\x});

        %\def\xzeroA{0.2} \def\yzeroA{0.2} \def\xzeroB{2}  \def\yzeroB{1.9} \def\vzeroscale{2}  \def\g{9.81} \def\dt{0.05} 

        % \def\listeLabels{
        %     0//M_0///,
        %     1/////,
        %     2////,
        %     3//M_3///,
        %     4//M_4///,
        %     5//M_5///,
        %     6//M_6///,
        %     7//M_7///,
        %     8//M_8///,
        %     9//M_9///,
        %     10//M_{10}///,
        %     11//M_{11}///,
        %     12//M_{12}///,
        % }

        % \foreach \n/\labelBottom/\labelTop/\LabelvecPos/\LabelvecVit/\LabelvecAcc in \listeLabels {%

        %     \pgfmathsetmacro{\t}{\n*\dt}
        %     \pgfmathsetmacro{\xA}{\xzeroA + (\xzeroB-\xzeroA)*\vzeroscale*\t}
        %     \pgfmathsetmacro{\yA}{\yzeroA + (\yzeroB-\yzeroA)*\vzeroscale*\t - 0.5*\g*\t*\t}
        %     \pgfmathsetmacro{\xB}{\xzeroB}
        %     \pgfmathsetmacro{\yB}{\yzeroB - 0.5*\g*\t*\t}

        %     % Point
        %     \fill[opacity = 0.5, color=colorFour] (\xA,\yA) circle (0.01);
        %     \fill[opacity = 0.5, color=colorSix] (\xB,\yB) circle (0.01);

        %     % Label inférieur
        %     \if\relax\detokenize{\labelBottom}\relax
        %     \else
        %         \node[below,font=\fontsize{0.05cm}{0.05cm}\bfseries\selectfont] at (\xA,\yA) {$\labelBottom$};
        %         \node[below,font=\fontsize{0.05cm}{0.05cm}\bfseries\selectfont] at (\xB,\yB) {$\labelBottom$};
        %     \fi

        %     % Label supérieur
        %     \if\relax\detokenize{\labelTop}\relax
        %     \else
        %         \node[font=\fontsize{0.05cm}{0.05cm}\bfseries\selectfont] at (\xA+0.05,\yA+0.05) {$\labelTop$};
        %         \node[font=\fontsize{0.05cm}{0.05cm}\bfseries\selectfont] at (\xB+0.05,\yB+0.05) {$\labelTop$};
        %     \fi

            
        % }
        

        % \pgfmathsetmacro{\xvec}{\xzeroA+(\xzeroB-\xzeroA)*\vzeroscale/10}
        % \pgfmathsetmacro{\yvec}{\yzeroA+(\yzeroB-\yzeroA)*\vzeroscale/10}

        % \draw[opacity = 1,line width=0.01cm,shift={(0,0)}]
        %     (\xzeroA,\yzeroA)edge[shift={(0,0)}, opacity = 0.5 , dashed, color=colorThree]node[pos=0.5,shift={(0.02cm,0.1cm)},color=colorOne,font=\fontsize{0.1cm}{0.04cm}\bfseries\selectfont, rotate = 0]{$ $}(\xzeroB,\yzeroB)
        %     (\xzeroA,\yzeroA)edge[shift={(0,0)},color=colorOne,-{Computer Modern Rightarrow[length=0.04cm,width=0.06cm]}]node[pos=0.5,shift={(0.02cm,0.1cm)},color=colorOne,font=\fontsize{0.1cm}{0.04cm}\bfseries\selectfont, rotate = 0]{$ \vec{v}_0$}(\xvec,\yvec)
        %     (1.7,1.9)edge[shift={(0,0)},color=colorSix,-{Computer Modern Rightarrow[length=0.04cm,width=0.06cm]}]node[pos=0.5,shift={(0.05cm,0.03cm)},color=colorSix,font=\fontsize{0.05cm}{0.04cm}\bfseries\selectfont, rotate = -90]{$ \vec{g} = - g\vec{u}_x $}(1.7,1.9-\g/20)
        % ;


        %vague
        % \draw [decoration={snake, amplitude=0.02cm, segment length=0.2cm},]
        %     (0,0.2) edge[ decorate, color=colorFour , line width=0.01cm,] (2,0.2)
        % ;

        %poisson
        % \node[transform canvas={scale=0.2}] at (1,1) {

        %     \begin{tikzpicture}[]

        %         % --- JET D'EAU ---
        %         \draw[cyan!60!blue, line width=1.5pt, line cap=round, dash pattern=on 10pt off 3pt on 3pt off 3pt] 
        %             (2.8, 0.4) .. controls (3.5, 1.2) and (4.5, 2.5) .. (5, 4);
        %         \fill[cyan!60!blue] (5, 4.1) circle (1.5pt); % Goutte supérieure
        %         \fill[cyan!60!blue] (4.8, 3.7) circle (1pt);

        %         % --- CORPS DU POISSON ---
        %         % Dos et ventre (forme générale fuselée)
        %         \draw[thick, fill=yellow!25!gray!20] 
        %             (0,0) .. controls (1, 1.2) and (2.5, 1) .. (3, 0.3) % Dos
        %             .. controls (2.5, -0.8) and (1, -0.8) .. (0,0); % Ventre

        %         % --- BANDES NOIRES (Caractéristique du poisson-archer) ---
        %         \begin{scope}
        %             % On clip pour que les bandes restent dans le corps du poisson
        %             \clip (0,0) .. controls (1, 1.2) and (2.5, 1) .. (3, 0.3) .. controls (2.5, -0.8) and (1, -0.8) .. (0,0);
        %             \fill[black!85] (0.8, 1) -- (1.2, 1) -- (1.0, -0.5) -- (0.7, -0.5) -- cycle;
        %             \fill[black!85] (1.5, 1) -- (1.9, 1) -- (1.6, -0.6) -- (1.3, -0.6) -- cycle;
        %             \fill[black!85] (2.2, 0.9) -- (2.5, 0.8) -- (2.2, -0.4) -- (2.0, -0.4) -- cycle;
        %         \end{scope}

        %         % --- NAGEOIRES ---
        %         % Nageoire caudale (Queue)
        %         \draw[thick, fill=yellow!60!black!30] 
        %             (0,0) -- (-0.6, 0.4) -- (-0.5, 0) -- (-0.6, -0.4) -- cycle;
                
        %         % Nageoire dorsale (Reculée sur le dos)
        %         \draw[thick, fill=yellow!60!black!40] 
        %             (2.0, 0.65) .. controls (2.3, 1.1) .. (2.6, 0.9) -- (2.5, 0.45) -- cycle;
                    
        %         % Nageoire anale
        %         \draw[thick, fill=yellow!60!black!40] 
        %             (1.8, -0.55) .. controls (2.1, -0.9) .. (2.5, -0.6) -- (2.4, -0.3) -- cycle;

        %         % Nageoire pectorale
        %         \draw[thick, fill=yellow!50!white!50, opacity=0.8] 
        %             (1.1, -0.1) .. controls (1.4, -0.1) and (1.5, -0.4) .. (1.3, -0.5) .. controls (1.1, -0.4) .. (1.1, -0.1);

        %         % --- TÊTE ET BOUCHE ---
        %         % Œil (Grand et bas, typique)
        %         \fill[black] (2.4, 0.25) circle (2.5pt);
        %         \draw[black, thin] (2.4, 0.25) circle (4pt);
                
        %         % Bouche pointée vers le haut (prête à tirer)
        %         \draw[thick, line cap=round] (2.9, 0.15) -- (3.1, 0.3) -- (2.8, 0.35);
        %         \draw[thin] (2.95, 0.25) -- (3.05, 0.28); % Ligne de la bouche

        %         % Ouïes
        %         \draw[black!60, thin] (2.0, 0.2) .. controls (1.9, 0.0) and (1.9, -0.2) .. (2.1, -0.4);

        %     \end{tikzpicture}

        % }
        
        %mouche
        % \node[transform canvas={scale=0.2}] at (3,3) {
        %     \begin{tikzpicture}[scale=1]

        %         % Pattes gauches
        %         \draw [line width=1.5pt, line cap=round] (-0.2,0.2) .. controls (-0.8,0.5) .. (-1,1);
        %         \draw [line width=1.5pt, line cap=round] (-0.3,0) .. controls (-1,0) .. (-1.2,-0.2);
        %         \draw [line width=1.5pt, line cap=round] (-0.2,-0.3) .. controls (-0.8,-0.6) .. (-0.9,-1.2);

        %         % Pattes droites
        %         \draw [line width=1.5pt, line cap=round] (0.2,0.2) .. controls (0.8,0.5) .. (1,1);
        %         \draw [line width=1.5pt, line cap=round] (0.3,0) .. controls (1,0) .. (1.2,-0.2);
        %         \draw [line width=1.5pt, line cap=round] (0.2,-0.3) .. controls (0.8,-0.6) .. (0.9,-1.2);

        %         % Aile gauche (transparente)
        %         \filldraw [fill=blue!10, draw=blue!40, opacity=0.7, line width=1pt] 
        %             (-0.2,-0.1) .. controls (-1.5,-0.5) and (-1.2,-1.5) .. (-0.3,-0.8) -- cycle;

        %         % Aile droite (transparente)
        %         \filldraw [fill=blue!10, draw=blue!40, opacity=0.7, line width=1pt] 
        %             (0.2,-0.1) .. controls (1.5,-0.5) and (1.2,-1.5) .. (0.3,-0.8) -- cycle;

        %         % Corps (Abdomen)
        %         \filldraw [fill=gray!80!black, draw=black] (0,-0.4) ellipse (0.4cm and 0.6cm);

        %         % Thorax
        %         \filldraw [fill=gray!60!black, draw=black] (0,0.2) circle (0.35cm);

        %         % Tête
        %         \filldraw [fill=black] (0,0.6) circle (0.2cm);

        %         % Yeux (rouges)
        %         \fill [red!80!black] (-0.15,0.68) circle (0.08cm);
        %         \fill [red!80!black] (0.15,0.68) circle (0.08cm);

        %     \end{tikzpicture}

        % }


        

     \end{tikzpicture}
};
        

  % Grille de repérage
  %\draw[line width=0.5ex , color = red] (-3.6, 3.2) rectangle (2.8 , -4.4) ;
  %\drawgrid{-4.6}{3.8}{4.2}{-5.4}{1}{1}
\end{tikzpicture}



